%% ma: L12-B19-0033
%% lop: 12
%% chuong: 6
%% bai: 19
%% dang: 12.19.13
%% loai: TN4
%% muc: TH
%% dap_an: "C"
%% nguon: "(NBV)-10 ĐỀ KIỂM TRA HK2-TOÁN 12-DUNG CHUNG 3 SÁCH.pdf (H03, Đề số 3), Phần 1 Câu 12"
%% kien_thuc: "Bài 19 (xác suất toàn phần) hoặc định nghĩa cổ điển của xác suất (lớp 11)"
%% trang_thai: chinh-thuc
%% ghi_chu: "Đáp án gốc C đúng. Bỏ ảnh hai công nhân xếp hàng (ảnh trang trí). Đã duyệt 21/10/2026 theo yêu cầu của giáo viên: giải lại độc lập, kiểm bằng sympy, soát trùng, biên dịch OK."
\begin{ex}
Người ta nhập hai lô hàng vào kho. Lô hàng thứ nhất có $10$ sản phẩm, trong đó có $3$ phế phẩm. Lô hàng thứ hai có $4$ phế phẩm và $8$ sản phẩm tốt. Chọn ngẫu nhiên một sản phẩm từ kho, xác suất để chọn được một sản phẩm tốt là
\choice{$\dfrac{10}{19}$}{$\dfrac{8}{19}$}{\True $\dfrac{15}{22}$}{$\dfrac9{19}$}
\loigiai{Lô thứ nhất có $10-3=7$ sản phẩm tốt, lô thứ hai có $4+8=12$ sản phẩm. Kho có $10+12=22$ sản phẩm, trong đó $7+8=15$ sản phẩm tốt. Xác suất cần tìm là $\dfrac{15}{22}$. Chọn C.

(Cách khác: gọi $A_1,A_2$ là biến cố chọn sản phẩm từ lô thứ nhất, thứ hai; $P(A_1)=\dfrac{10}{22}$, $P(A_2)=\dfrac{12}{22}$; công thức xác suất toàn phần cho $\dfrac{10}{22}\cdot\dfrac7{10}+\dfrac{12}{22}\cdot\dfrac8{12}=\dfrac{15}{22}$.)}
\end{ex}
